\documentclass{article}

\usepackage{amsmath,amssymb}
\usepackage{xurl}
\usepackage[hidelinks]{hyperref}
\setlength{\emergencystretch}{2em}

% Shared commands for notes in docs/paper-gaps/.

\DeclareUnicodeCharacter{2080}{\ifmmode{}_{0}\else\textsubscript{0}\fi}
\DeclareUnicodeCharacter{2081}{\ifmmode{}_{1}\else\textsubscript{1}\fi}
\DeclareUnicodeCharacter{2082}{\ifmmode{}_{2}\else\textsubscript{2}\fi}
\DeclareUnicodeCharacter{2099}{\ifmmode{}_{n}\else\textsubscript{n}\fi}

\newcommand{\Lean}{\textsc{Lean}}
\ExplSyntaxOn
\NewDocumentCommand{\leanid}{m}{
  \ifmmode
    \textsf{\detokenize{#1}}
  \else
    \group_begin:
    \urlstyle{sf}
    \tl_set:Nn \l_tmpa_tl {#1}
    \tl_replace_all:Nnn \l_tmpa_tl {\_} {_}
    \exp_args:NV \path \l_tmpa_tl
    \group_end:
  \fi
}
\ExplSyntaxOff
\providecommand{\tr}{\operatorname{tr}}
\newcommand{\tnleanIssueBase}{https://github.com/LionSR/TNLean/issues/}
\newcommand{\tnleanPrBase}{https://github.com/LionSR/TNLean/pull/}
\newcommand{\tnleanDocBase}{https://sirui-lu.com/TNLean/docs/}
\newcommand{\ghissue}[1]{\href{\tnleanIssueBase#1}{GitHub issue \##1}}
\newcommand{\ghpr}[1]{\href{\tnleanPrBase#1}{PR \##1}}
\newcommand{\doclink}[1]{\href{\tnleanDocBase#1.html}{\path{#1}}}

% Dirac notation and the identity operator, as in blueprint/src/macros/common.tex.
% \providecommand keeps note-local definitions authoritative.
\usepackage{dsfont}
\providecommand{\ket}[1]{|#1\rangle}
\providecommand{\bra}[1]{\langle #1|}
\providecommand{\braket}[2]{\langle #1 | #2 \rangle}
\providecommand{\ketbra}[2]{|#1\rangle\!\langle #2|}
\providecommand{\Id}{\mathds{1}}
\providecommand{\one}{\mathds{1}}
\providecommand{\C}{\mathbb{C}}
\providecommand{\MN}[1]{M_{#1}(\C)}
\providecommand{\MD}{M_D(\C)}

% External referencing for paper sources.  The build helper writes sanitized
% auxiliary files under build/paper-aux/ and excludes duplicated source labels.
\usepackage{xr}

% Import a generated paper auxiliary file with a caller-supplied label prefix.
% The candidate paths support builds from docs/paper-gaps/, the repository root,
% and build/paper-aux/ itself.
\newcommand{\importpaperaux}[2]{%
  \IfFileExists{../../build/paper-aux/#2.aux}{%
    \externaldocument[#1]{../../build/paper-aux/#2}%
  }{%
    \IfFileExists{build/paper-aux/#2.aux}{%
      \externaldocument[#1]{build/paper-aux/#2}%
    }{%
      \IfFileExists{#2.aux}{%
        \externaldocument[#1]{#2}%
      }{}%
    }%
  }%
}

% General external reference macros
\newcommand{\paperref}[2]{\ref{#1-#2}}
\newcommand{\papereqref}[2]{(\ref{#1-#2})}

% CPSV16 source (1606.00608 / MPDO-22-12-17-2.tex) external document and shortcuts
\importpaperaux{cpsv16-}{1606.00608/MPDO-22-12-17-2-tnlean}
\newcommand{\cpsvref}[1]{\paperref{cpsv16}{#1}}
\newcommand{\cpsveqref}[1]{\papereqref{cpsv16}{#1}}

% Verdict marker: \gapnote{<kind>}{<status>}. Kind is the policy
% classification (clarification, local-correction, scope-restriction,
% unfaithful, false-source, open-gap); status is open, wip (elimination
% underway), resolved, or historical. Machine-read by the site index;
% prints a verdict line.
\newcommand{\gapnote}[2]{\par\noindent\textbf{Verdict:} #1 (#2).\par}


\title{The Two-Dimensional Overlapping-Window Strengthening of the
Normal PEPS Fundamental Theorem}
\author{}
\date{2026-06-13}

\begin{document}
\maketitle
\gapnote{open-gap}{open}

\section{The source statement}

The final unformalized statement of \cite{Molnar2018NormalPEPS} is the
two-dimensional strengthening of the torus Fundamental Theorem announced
at the end of the overlapping-window
section.\footnote{\path{Papers/1804.04964/paper_normal.tex:2296--2318}.}
Ref.~\cite{Molnar2018NormalPEPS} states the corollary as follows:

\begin{quote}
  Let $A$ and $B$ be two normal 2D PEPS tensors such that every
  $L \times K$ region is injective. Suppose they generate the same state
  on some region $n\times m$ with $n \geq 2L+1$ and $m\geq 2K+1$. Then
  $A$ and $B$ are related to each other with a gauge:
  \begin{align}
    B^i
      &= \lambda \cdot (X \otimes Y)\, A^i\,
         (X^{-1} \otimes Y^{-1}),
    \label{eq:peps_normal_ft_2d_overlap_source_gauge_relation}
  \end{align}
  with $\lambda^{n\cdot m} = 1$ and $X,Y$ invertible matrices. Moreover
  these matrices $X,Y$ are unique up to a multiplicative constant.
\end{quote}

The relation (\ref{eq:peps_normal_ft_2d_overlap_source_gauge_relation})
abbreviates the diagram in Ref.~\cite{Molnar2018NormalPEPS}: $X$ acts on the
two horizontal virtual legs of the tensor and $Y$ on the two vertical legs,
with inverses on the outgoing sides. Its conclusion is identical to that of
Theorem~3 of Ref.~\cite{Molnar2018NormalPEPS};\footnote{Theorem~3 is at
\path{Papers/1804.04964/paper_normal.tex:1451--1472}; its torus form is
formalized as
\leanid{TNLean.PEPS.fundamentalTheorem_normalTorusPEPS_unconditional} in
\path{TNLean/PEPS/TorusFundamentalTheorem2.lean}.} only the hypotheses
differ. Theorem~3 of Ref.~\cite{Molnar2018NormalPEPS} assumes that every
$2\times3$ \emph{and} every $3\times2$ region is injective and that
$n, m \geq 7$. The corollary of Ref.~\cite{Molnar2018NormalPEPS} assumes
injectivity of a single rectangle shape $L \times K$ --- one
orientation only --- and the sizes $n \geq 2L+1$, $m \geq 2K+1$. For
$L = K = 2$ this reads: injectivity of every $2\times2$ region and sizes
$n, m \geq 5$, strictly weaker hypotheses than the formalized theorem in
both the region shapes and the sizes. The strengthening is therefore
real even in the smallest case, and grows with $L, K$ (for $L=2$,
$K=3$ it gives $5 \times 7$ in place of $7 \times 7$).

\section{The source proof sketch}

Ref.~\cite{Molnar2018NormalPEPS} gives only a proof sketch.\footnote{\path{Papers/1804.04964/paper_normal.tex:2320--2445}.}
Its full text, with the diagrams described in words:

\begin{quote}
  We only need to prove a statement similar to Lemma~5. For that, notice
  that a virtual operation on a given bond can be interpreted as a
  physical operation on any of the following four regions (in the case
  of $K=L=2$): [four $2\times2$ windows around a horizontal edge: the
  window right of the edge, then sliding left across the edge in two
  steps, then sliding down one step]. We need to prove that conversely,
  any four physical operators on the above regions that transforms the
  PEPS into the same state means that the transformation can equally be
  done with a virtual operation on the highlighted edge. The system size
  required to compare any two consecutive regions is only $5\times 5$
  (and in general, $(2L+1)\times (2K+1)$). Therefore, similar to
  Lemma~5, [the four operators applied to the patch around the edge
  agree] with open boundaries. Compare now the first and the last
  expression in the above equation. One can add two-two tensors in the
  upper left and lower right corner and invert the resulting regions,
  leading to [the equality restricted to the staircase-shaped union of
  the first and last windows]. This results in the desired virtual
  operation on the highlighted edge. The rest of the proof is the same
  as the proof of Theorem~3.
\end{quote}

Lemma~5 of Ref.~\cite{Molnar2018NormalPEPS} is the one-dimensional
overlapping-window bond-operator extraction,\footnote{\path{Papers/1804.04964/paper_normal.tex:2045--2255},
formalized (site-dependent form) as
\leanid{MPSChainTensor.overlapWindow_exists_bondOperator} in
\path{TNLean/PEPS/CycleMPSChainOverlapWindow.lean}, with the closed-chain
corollary \leanid{TNLean.PEPS.fundamentalTheorem_normalMPSChain_of_overlap}
in \path{TNLean/PEPS/CycleMPSChainOverlapCapstone.lean}.} which delivers
the one-dimensional Fundamental Theorem at $n \geq 2L+1$ in place of the
three-block bound $n \geq 3L$.

\subsection{What the sketch asserts, and what remains to be extracted}

The sketch in Ref.~\cite{Molnar2018NormalPEPS} is explicitly labelled
``Sketch of proof''\footnote{\path{Papers/1804.04964/paper_normal.tex:2320}.} and its
single load-bearing claim is asserted by analogy to Lemma~5 rather than
derived. The claim in Ref.~\cite{Molnar2018NormalPEPS} is the converse
correspondence between virtual and physical operations: ``a virtual
operation on a given bond can be
interpreted as a physical operation on any of the [\,$L+K$\,] regions,'' and
conversely ``any [\,$L+K$\,] physical operators on the above regions that
transforms the PEPS into the same state means that the transformation can
equally be done with a virtual operation on the highlighted
edge.''\footnote{\path{Papers/1804.04964/paper_normal.tex:2321} and
\path{Papers/1804.04964/paper_normal.tex:2368}.} The formalization separates
the two directions of this correspondence.  The forward direction ---
realizing a fixed virtual operation $X$ by a physical operator $O_j(X)$ on
every staircase window --- is now complete: the open-boundary identity
$O_j(X)T_{B,W_j}(\mu)=C_j(X)(\mu)$ is proved for every virtual boundary
configuration~$\mu$.  The converse, which must compare the two tensor
families and recover the algebra homomorphism carried by $X$, is the
remaining extraction problem.

Ref.~\cite{Molnar2018NormalPEPS} supplies this converse by the reduction ``we only need to prove a
statement similar to Lemma~5''\footnote{\path{Papers/1804.04964/paper_normal.tex:2321}.}
and ``therefore, similar to Lemma~5, [the four operators on the patch
agree] with open boundaries,''\footnote{\path{Papers/1804.04964/paper_normal.tex:2368}.}
followed by the corner completion ``one can add two-two tensors in the upper
left and lower right corner and invert the resulting
regions.''\footnote{\path{Papers/1804.04964/paper_normal.tex:2415}.} Two of
these three steps transfer. The consecutive-window agreement with open
boundaries, obtained by inverting the complement of each consecutive-window
union, is rigorous and needs only the sizes $n \geq 2L+1$, $m \geq 2K+1$ (the
sketch's ``$5\times5$'' / ``$(2L+1)\times(2K+1)$''); the corner completion is
rigorous as a cancellation of one injective $L \times K$ rectangle. Both are
formalized; see the filled-in derivation below.

The step that does \emph{not} transfer is the final extraction of the single
matrix $X$ on the edge, which the sketch obtains by analogy to the closing
move of Lemma~5. In Lemma~5 of Ref.~\cite{Molnar2018NormalPEPS}, that move compares the
two ends of the chain, ``similar to'' the injective read-off,\footnote{\path{Papers/1804.04964/paper_normal.tex:2213}, invoking
\path{Papers/1804.04964/paper_normal.tex:377}, source equation
\path{eq:inj_O->X_argument} of Ref.~\cite{Molnar2018NormalPEPS}.} which there yields the bond matrix $X$ and,
by composition, the algebra
homomorphisms $O_1 \mapsto X$ and
$O_3^T \mapsto X$.\footnote{\path{Papers/1804.04964/paper_normal.tex:2129} and
\path{Papers/1804.04964/paper_normal.tex:2252} of
Ref.~\cite{Molnar2018NormalPEPS}.} That read-off is valid in
one dimension for a structural reason specific to a one-dimensional window:
a length-$L$ window has exactly two virtual legs, a left bond and a right
bond, both of dimension $D$, so inverting the injective window tensors
produces a genuine $D \times D$ \emph{matrix} $W$ on the bond,\footnote{The
inverted-window matrix is the $W$ of
\path{Papers/1804.04964/paper_normal.tex:401}--\path{403}, an
$D_{23} \times D_{23}$ object on the single bond.} and window injectivity is
precisely the statement that these matrices span the full bond algebra, which
is what makes the homomorphism extraction go through.

In two dimensions an end window used for the bond read-off no longer has this
structure.  An end $L \times K$ window touches the reference edge $e$ with \emph{one}
endpoint, so the leg of $e$ inside an end window is a single
$\mathrm{Fin}(\mathrm{bondDim}\,e)$ index, not a pair; reading the window
block at fixed external legs gives a \emph{vector} on
$\mathrm{Fin}(\mathrm{bondDim}\,e)$, never a square matrix on the bond, and
the rest of the window's perimeter is open. The two indices of the matrix $X$
the corollary wants are the $e$-legs of the two diagonally offset end
windows, one from each; no single window supplies a square-matrix family, and
folding a window's external perimeter into a second
$\mathrm{Fin}(\mathrm{bondDim}\,e)$ index would require
$\prod_{\text{external legs}}\mathrm{bondDim} = \mathrm{bondDim}\,e$, false in
general.  Thus a direct port of the one-dimensional single-window matrix-span
proof is not available.  This rules out that particular proof path; it does
not show that the paper's overlapping-window comparison is false.  In
particular, it must not be conflated with the forward virtual-operation
family, which keeps the same distinguished edge $e$ throughout and requires
no such span.

\medskip
\noindent\textbf{Formalization verdict.} The published argument is explicitly
a sketch.  The consecutive-window comparison, corner completion, and forward
open-boundary physical realization of one virtual operation on all windows are
now formalized.  What
remains is the cross-tensor converse and the homomorphism extraction
corresponding to Lemma~5's maps $O_1\mapsto X$ and
$O_3^{\mathsf T}\mapsto X$.  The failed single-window span above is an
invented auxiliary route, not a statement made or required by the paper.

\section{Consolidated account}

This section is the self-contained summary of the investigation. The
remaining sections record, in full, the machinery that was built and the
machine-checked refutations of each closing route; this section states the
conclusions a reader needs without that detail.

\paragraph{1. The statement and its faithful Lean hypotheses.}
The corollary of Ref.~\cite{Molnar2018NormalPEPS}\footnote{\path{Papers/1804.04964/paper_normal.tex:2297--2318}.}
asserts that two normal $2$D PEPS tensors $A$, $B$ with every $L \times K$
region injective, generating the same state on an $n \times m$ region with
$n \geq 2L+1$, $m \geq 2K+1$, are related by
$B^i = \lambda\,(X \otimes Y)\,A^i\,(X^{-1} \otimes Y^{-1})$ with
$\lambda^{nm} = 1$ and $X, Y$ unique up to a constant. ``Normal $\dots$ such
that every $L \times K$ region is injective'' is the source's own normality
condition in Ref.~\cite{Molnar2018NormalPEPS}\footnote{\path{Papers/1804.04964/paper_normal.tex:1318}: ``We call
a PEPS \emph{normal}, if blocking tensors in certain regions results in
injective tensors.''} instantiated at the $L \times K$ blocking; it is not a
second, stronger hypothesis. The faithful Lean hypothesis bundle is
\leanid{TNLean.PEPS.NormalTorusArcWindowInjectivityHypotheses}
(\path{TNLean/PEPS/TorusWindowComplement.lean}), a single field
\leanid{arcWindow_injective}: every cyclic $L \times K$ window is injective in
the sense of \leanid{TNLean.PEPS.RegionBlockedTensorInjective}. This neither
exceeds nor falls short of the source's hypothesis.

\paragraph{2. What is formalized.}
All surrounding machinery is built and lands at the corollary's minimal size:
the seam-aware window geometry
(\path{TNLean/PEPS/TorusWindowRegion.lean},
\path{TNLean/PEPS/TorusWindowComplement.lean}); the deformed-window state
vocabulary and the consecutive-window comparison by complement inversion
(\path{TNLean/PEPS/TorusDeformedWindow.lean}); the chaining, corner
completion, and shared-corner cancellation that produce the
\emph{open-boundary} end-pair equality
\leanid{TNLean.PEPS.staircasePair_insert_eq_open}
(\path{TNLean/PEPS/TorusWindowChain2.lean} through
\path{TorusWindowChain6.lean}) without ever inverting the non-injective
$\mathrm{univ} \setminus S$; the single-bond peeling of that equality onto the
reference edge
(\leanid{TNLean.PEPS.regionInsertedCoeff_endWindows_eq_of_staircase},
\path{TNLean/PEPS/TorusWindowPeeling/SingleBond.lean}); the per-edge witness packaging and
the Skolem--Noether read-off
(\path{TNLean/PEPS/TorusWindowWitness.lean},
\path{TNLean/PEPS/TorusWindowMult.lean}); the per-step bond-insert transport
\leanid{TNLean.PEPS.horizontalStaircaseConsecutiveWindow_bondTransport_extend_eq}
and the now-uniform interior-bond rescaling
\leanid{TNLean.PEPS.regionInteriorBondProd_staircaseWindow_eq}
(\path{TNLean/PEPS/TorusWindowBondTransport.lean},
\path{TorusWindowBondUniform.lean}); the interior-bond insert and the complete
common-state family
(\leanid{TNLean.PEPS.interiorBondInsertedRegionInsert} and
\leanid{TNLean.PEPS.staircaseVirtualOperationInsert}), together with its
open-boundary physical operators
(\leanid{TNLean.PEPS.physicalOpOfRegionInsert} and
\leanid{TNLean.PEPS.staircaseVirtualOperationPhysicalOp\_realizes}), and its
end-window coefficient identity
\leanid{TNLean.PEPS.regionInsertedCoeff_endWindows_eq_staircaseVirtualOperation}
(\path{TNLean/PEPS/RegionBlock/InteriorBondInsertion.lean} and
\path{TNLean/PEPS/TorusWindowVirtualOperationFamily.lean}); and the back-half
assembly that re-anchors
the torus Theorem~3 layers to the $(L+1) \times (K+1)$ comparison pair. Every
inverted region is a union of one-orientation $L \times K$ window translates,
host-decomposable at $n \geq 2L+1$, $m \geq 2K+1$. The development is
minimal-size-ready through the single-tensor virtual-operation family.  The
remaining gate is the cross-tensor transfer and its multiplicativity.

\paragraph{3. The remaining cross-tensor extraction.}
The same-tensor operation family is no longer missing.  For a fixed matrix
$X$ on $e$, the first window carries $X^{\mathsf T}$, the windows
$W_1,\ldots,W_{L-1}$ carry the interior-bond insert of $X$, and the remaining
windows carry the boundary insert of $X$.  Window injectivity defines
$O_j(X)$ and proves $O_j(X)T_{B,W_j}(\mu)=C_j(X)(\mu)$ before complementary
contraction; all these inserts also have common closed state
$P_B(W_0)E_B(e,\cdot,X)$.  The outstanding input is instead the cross-tensor
coefficient transfer between $A$ and $B$ and its multiplicativity
\leanid{hmul}, the two-dimensional counterpart of Lemma~5's algebra
homomorphisms.  Existing single-window and coarse-three-site attempts do not
derive that input at the minimal size; a source-faithful extraction must use
the actual overlapping-window comparison rather than assume a matrix span.

\paragraph{4. Auxiliary closing routes ruled out at the minimal size.}
The following findings rule out particular substitutes for the paper's
overlapping-window argument.  They do not constitute a proof that the source
corollary is false or that no source-faithful proof exists.
\begin{itemize}
  \item \emph{Single-vertex injectivity is unavailable.} Normality is by
    definition \emph{weaker} than single-vertex injectivity: a normal tensor
    need not be injective at one site but becomes injective after blocking.
    So normality cannot supply
    \leanid{TNLean.PEPS.IsVertexInjective}, and the endpoint span it would
    feed
    (\leanid{TNLean.PEPS.span_stateOpenCoeff_eq_top}) is not derivable from
    the hypotheses.
  \item \emph{The one-dimensional port is absent.} The one-dimensional
    extraction
    \leanid{MPSChainTensor.overlapWindow_exists_bondOperator} succeeds from
    window injectivity alone because a length-$L$ one-dimensional window
    product is a square matrix on its bond. A two-dimensional $L \times K$
    window touches the bond with one endpoint and gives a vector family, never
    a square-matrix span (Part~1).
  \item \emph{The end-pair coarse frame is non-injective.} Taking the red
    block to be one end window forces the single-crossing partner to be the
    opposite end window, so the third region of the frame is
    $\mathrm{univ} \setminus S$, which is not injective at the minimal size:
    the end pair occupies all $2L$ columns of one row and the complement band
    there is below $L$ columns wide.
  \item \emph{No per-bond coarse frame exists, for any single-crossing window
    pair.} The only landed source of complement injectivity at the minimal
    size needs $\mathrm{red} \cup \mathrm{blue}$ to be a single cyclic
    rectangle, but two $L \times K$ windows cross at a single edge only when
    they are the diagonally offset pair, whose union is a staircase, not a
    rectangle
    (\leanid{TNLean.PEPS.horizontalStaircasePair_ne_arcRectangle}). This is a
    joint unsatisfiability, settled combinatorially in
    \path{TNLean/PEPS/TorusWindowSingleCrossingObstruction.lean}.
  \item \emph{Translation invariance supplies the common normalization.}
    The corollary's torus translation invariance equalizes the non-boundary
    bond products of all staircase windows
    (\leanid{TNLean.PEPS.regionInteriorBondProd_staircaseWindow_eq}).  Together
    with the interior-bond insert, this completes the common-state family; it
    does not by itself compare the two tensor families.
  \item \emph{The row-and-column one-dimensional reduction fails.} Treating
    each row as one coarse site and applying the one-dimensional theorem
    row-wise and column-wise cannot reach the per-edge conclusion: a
    contracted row is invariant under conjugating its bond family by an
    arbitrary invertible matrix, so the reduction cannot distinguish a
    per-edge-gauged row from a non-product conjugate that has the same state
    and injectivity but no per-edge factorization. This is a machine-checked
    refutation
    (\leanid{TNLean.PEPS.RowColumnReductionObstruction.rowCutGaugeFactorizes_false},
    \path{TNLean/PEPS/TorusRowColumnReductionObstruction.lean}).
\end{itemize}

\paragraph{5. Reading the source faithfully.}
Ref.~\cite{Molnar2018NormalPEPS} keeps one highlighted edge fixed and compares physical operations
on overlapping regions; it does not introduce a chain of different bonds or
a single-window square-matrix span.  The forward correspondence, its
open-boundary physical operators, and its transpose convention now follow
this description literally.  The remaining
formal work is to expand the sketch's converse comparison far enough to obtain
the cross-tensor homomorphism, using Lemma~5's $O_1$ and
$O_3^{\mathsf T}$ notation and composition argument.

\paragraph{6. Status.}
The corollary remains \emph{unformalized}: no blueprint entry claims it, it
carries the \texttt{notready} marker, and no \texttt{leanok} marks it.  The
same-tensor family requested in issue~\#2780 is complete; the remaining work
is the cross-tensor transfer and homomorphism extraction tracked separately.

\section{The filled-in derivation}

This section records the complete argument the sketch compresses, in the
notation of the formal torus development: the lattice is the discrete
$n \times m$ torus, regions are sets of vertices, $T_R$ denotes the
tensor of the network blocked over a region $R$, and a region is
injective when $T_R$ is injective as a map from its virtual boundary
indices to its physical indices.

\subsection{The window family around an edge}

Fix a horizontal edge $e$ with endpoints $u = (x,y)$ and $w = (x+1,y)$.
Define $L+K$ windows, each a contiguous $L \times K$ rectangle:
\begin{align}
  W_j
    &= [x+1-j,\, x+L-j] \times [y,\, y+K-1],
    & j &= 0, \dots, L,
  \label{eq:peps_normal_ft_2d_overlap_horizontal_window_family}\\
  W_{L+i}
    &= [x-L+1,\, x] \times [y-i,\, y+K-1-i],
    & i &= 1, \dots, K-1.
  \label{eq:peps_normal_ft_2d_overlap_vertical_window_family}
\end{align}
The windows in (\ref{eq:peps_normal_ft_2d_overlap_horizontal_window_family})
slide left across the edge one column at a time; those in
(\ref{eq:peps_normal_ft_2d_overlap_vertical_window_family}) descend one row
at a time. The first window $W_0$
contains only the right endpoint $w$ (at its lower-left corner), the
windows $W_1, \dots, W_{L-1}$ contain both endpoints, and the windows
$W_L, \dots, W_{L+K-1}$ contain only the left endpoint $u$. For
$K = L = 2$ these are exactly the four regions of the sketch. Around a
vertical edge the construction is transposed in the roles of the two
axes with the same $L \times K$ window shape: $K+1$ vertical slides
followed by $L-1$ horizontal ones.

A virtual operation $X$ on the bond $e$ is represented on every window by
keeping the same distinguished edge fixed.  The first window, which contains
only the ordered right endpoint, carries the boundary insert of
$X^{\mathsf T}$; the windows $W_1,\dots,W_{L-1}$ carry the interior-bond
insert obtained by adjoining their genuine tensors to a one-endpoint insert;
and the remaining windows carry the boundary insert of $X$.  Translation
invariance identifies the non-boundary bond products, so all these inserts
produce the common state $P_A(W_0)E_A(e,\cdot,X)$.  This is formalized in
\path{TNLean/PEPS/TorusWindowVirtualOperationFamily.lean}.

The two-dimensional analogue of Lemma~5 also has a converse. Suppose each
window $W_j$ carries a deformed tensor $C_j$ (an arbitrary tensor with
the boundary structure of $W_j$) and all the deformed states agree:
replacing $T_{W_j}$ by $C_j$ in the torus contraction yields one common
state for $j = 0, \dots, L+K-1$. Then there is a single matrix $X$ on the bond $e$ satisfying
\begin{align}
  C_0
    &= T_{W_0} \cdot X_e,
  \label{eq:peps_normal_ft_2d_overlap_first_window_factorization}\\
  C_{L+K-1}
    &= X_e \cdot T_{W_{L+K-1}}.
  \label{eq:peps_normal_ft_2d_overlap_last_window_factorization}
\end{align}
Here $X_e$ in (\ref{eq:peps_normal_ft_2d_overlap_first_window_factorization})
and (\ref{eq:peps_normal_ft_2d_overlap_last_window_factorization}) denotes
the insertion of $X$ on the open leg of $e$, and
$X$ is unique.

\subsection{Step 1: comparing consecutive windows}

For consecutive windows the union $U_j = W_j \cup W_{j+1}$ is an
$(L+1) \times K$ rectangle (horizontal slides) or an $L \times (K+1)$
rectangle (vertical slides). The torus complement of $U_j$ is injective:
it decomposes into a full-height band of $L$ (respectively $L+1$)
columns and a leftover rectangle of size $(L+1)\times(K+1)$
(respectively $L \times K$), and each piece is a union of translated
$L \times K$ rectangles. The decomposition needs exactly $n \geq 2L+1$
and $m \geq 2K+1$; this is the only step of the whole argument where
the torus size enters, matching the bound
``$(2L+1)\times(2K+1)$'' in Ref.~\cite{Molnar2018NormalPEPS}. The bands may wrap the torus seam, so the
covering rectangles are translates of coordinate rectangles;
injectivity transports along translations.

Inverting the complement strips the state equality for $j$ and $j+1$ to
the open-boundary equality
\begin{align}
  C_j \,\sqcup\, T_{W_{j+1} \setminus W_j}
    &= T_{W_j \setminus W_{j+1}} \,\sqcup\, C_{j+1}
    &&\text{as tensors on } U_j.
  \label{eq:peps_normal_ft_2d_overlap_consecutive_window_open_equality}
\end{align}
Equation~(\ref{eq:peps_normal_ft_2d_overlap_consecutive_window_open_equality})
is the two-dimensional form of source equations \path{eq:normal_act_123} and
\path{eq:normal_act_234} in the proof of Lemma~5 of
Ref.~\cite{Molnar2018NormalPEPS}.\footnote{The one-dimensional inversions
are at \path{Papers/1804.04964/paper_normal.tex:2133--2179}.}

\subsection{Step 2: chaining to the patch}

Let $P = \bigcup_j W_j$, the staircase-shaped patch. Contracting each
instance of
(\ref{eq:peps_normal_ft_2d_overlap_consecutive_window_open_equality}) with
the tensors of $P \setminus U_j$
extends it to $P$, and transitivity chains the $L+K-1$ equalities into
the end-window comparison
\begin{align}
  C_0 \,\sqcup\, T_{P \setminus W_0}
    &= T_{P \setminus W_{L+K-1}} \,\sqcup\, C_{L+K-1}
    &&\text{on } P.
  \label{eq:peps_normal_ft_2d_overlap_patch_end_window_comparison}
\end{align}
Equation~(\ref{eq:peps_normal_ft_2d_overlap_patch_end_window_comparison})
has open boundary. It is the analogue of the insertion of $A_4$ and $A_1$ in
the proof of Lemma~5 of Ref.~\cite{Molnar2018NormalPEPS}.

\subsection{Step 3: stripping to the staircase pair}

Let $S = W_0 \sqcup W_{L+K-1}$, the disjoint union of the two end
windows. The difference $P \setminus S$ is the single block
$[x-L+1, x] \times [y+1, y+K-1]$ of size $L \times (K-1)$ (the sketch's
``upper left corner''; the sketch presents the comparison on the full
bounding box, where the lower-right corner block appears as well, and
removes both). An $L \times (K-1)$ block is not injective, so it cannot
be inverted directly. Instead, complete it: contract both sides of the
equality (\ref{eq:peps_normal_ft_2d_overlap_patch_end_window_comparison})
with one added row of network tensors above the block,
forming the $L \times K$ rectangle $[x-L+1,x] \times [y+1, y+K]$. The
completed rectangle is injective, and inverting it removes both the
added row and the corner block. This is the sketch's ``add two-two
tensors in the \dots corner and invert the resulting regions''. The
result is the open-boundary equality
\begin{align}
  C_0 \,\sqcup\, T_{W_{L+K-1}}
    &= T_{W_0} \,\sqcup\, C_{L+K-1}
    &&\text{on } S.
  \label{eq:peps_normal_ft_2d_overlap_staircase_pair_open_equality}
\end{align}

The derivation of
(\ref{eq:peps_normal_ft_2d_overlap_staircase_pair_open_equality}) cancels a
\emph{shared} injective block. Both
sides of (\ref{eq:peps_normal_ft_2d_overlap_patch_end_window_comparison}) carry the genuine network block on the
completed rectangle $Q = [x-L+1,x] \times [y+1,y+K]$ once the added row
is contracted in; they differ only on $S$. Injectivity of $Q$ then
cancels the common $Q$-block from both sides and leaves (\ref{eq:peps_normal_ft_2d_overlap_staircase_pair_open_equality}). This is a different inversion from the consecutive-window step of
Step~1, where the inverted region is the entire torus complement of the
compared region and the two sides differ across that whole complement.

\medskip
\noindent\textbf{The complement of the end pair is not injective at the
minimal size.} A natural but incorrect reading of the corner completion
is that it makes the torus complement $\mathrm{univ} \setminus S$
injective, so that the consecutive-window engine of Step~1 (which inverts
the full complement of the compared region) applies with the compared
region taken to be $S$ itself. This fails at the corollary's minimal
size. Read the cyclic columns and rows from the lower-left corner of the
left end window. The left window occupies columns $[0,L)$ at rows
$[0,K)$ and the right window occupies columns $[L,2L)$ at rows
$[K-1,2K-1)$, so at the single row $K-1$ the pair occupies \emph{all}
columns $[0,2L)$. The complement of $S$ in that row is the column band
$[2L, \mathrm{width})$, of width $\mathrm{width} - 2L$, which is exactly
$1$ at $\mathrm{width} = 2L+1$ and in general below $L$ whenever
$\mathrm{width} < 3L$. Every contiguous $L \times K$ window (or wider) that
contains a complement vertex at row $K-1$ must include a column in
$[0,2L)$ at that row, and every such column lies in $S$. So no injective
$L \times K$ translate fits through row $K-1$ of the complement, and
$\mathrm{univ} \setminus S$ is not a union of injective window translates;
its blocked tensor is not injective under the one-orientation window
hypotheses alone at this size. (For $\mathrm{width} \geq 3L$ and
$\mathrm{height} \geq 3K$ the full-width and full-height gaps both reach
$L$ and $K$, and $\mathrm{univ} \setminus S$ does decompose; the corollary
needs the minimal size where it does not.) The same obstruction rules out
the variant $\mathrm{univ} \setminus S = (\mathrm{univ} \setminus P) \cup
\mathrm{corner}$: the patch $P$ also occupies all columns $[0,2L)$ at row
$K-1$, so $\mathrm{univ} \setminus P$ has the same sub-$L$ band there.

The faithful step is therefore the shared-block cancellation that yields
(\ref{eq:peps_normal_ft_2d_overlap_staircase_pair_open_equality}), which
needs only injectivity of the completed rectangle $Q$ (an $L \times K$
window, injective by hypothesis) and never asserts injectivity of
$\mathrm{univ} \setminus S$. The Lean stripping is the open-boundary
shared-corner cancellation \leanid{staircasePair_insert_eq_open}: it cancels
the shared injective completed rectangle $Q$ from
(\ref{eq:peps_normal_ft_2d_overlap_patch_end_window_comparison}) and never
inverts $\mathrm{univ} \setminus S$. An earlier closed-torus stripping
discharged Step~3 with the full-complement engine
\leanid{deformedRegionStateAssembled_insert_eq_of_complementInjective} and so
carried \leanid{RegionBlockedTensorInjective}~$B\;(\mathrm{univ} \setminus
\text{\leanid{horizontalStaircaseEndPair}}\;s\,L\,K)$ as an explicit
hypothesis not derivable at the minimal size; it has been removed in favour of
the shared-injective-block cancellation engine for $Q$, which inverts the
common completed-corner block rather than the whole complement.

\subsection{Step 4: the single-crossing extraction}

The two end windows are the diagonally offset $L \times K$ rectangles
\begin{align}
  W_0
    &= [x+1, x+L] \times [y, y+K-1],
  \label{eq:peps_normal_ft_2d_overlap_right_end_window}\\
  W_{L+K-1}
    &= [x-L+1, x] \times [y-K+1, y].
  \label{eq:peps_normal_ft_2d_overlap_left_end_window}
\end{align}
The column ranges in (\ref{eq:peps_normal_ft_2d_overlap_right_end_window})
and (\ref{eq:peps_normal_ft_2d_overlap_left_end_window}) are disjoint, and
their row ranges intersect exactly
in $\{y\}$, so the \emph{only} lattice edge joining them is $e$ itself:
a horizontally adjacent pair must sit in columns $x$ and $x+1$ at a
common row in $[y, y+K-1] \cap [y-K+1, y] = \{y\}$, and a vertically
adjacent pair would need a common column. Equation
(\ref{eq:peps_normal_ft_2d_overlap_staircase_pair_open_equality}) therefore
pairs two injective tensors across the single bond
$e$, and the standard two-sided extraction --- combine one family by coefficients
representing the identity, exactly as in source equation
\path{eq:inj_O->X_argument} of Ref.~\cite{Molnar2018NormalPEPS}\footnote{The
source equation is at \path{Papers/1804.04964/paper_normal.tex:377}; the
matrix-level form is
\leanid{MPSChainTensor.exists_bondOperator_of_intertwine_span}.} ---
produces the unique matrix $X$ on $e$ with the factorizations
(\ref{eq:peps_normal_ft_2d_overlap_first_window_factorization}) and
(\ref{eq:peps_normal_ft_2d_overlap_last_window_factorization}). This
single-crossing geometry is the reason the extracted operator lives on
one edge: no multi-edge cut is ever inverted across.

\medskip
\noindent\textbf{The single-crossing geometry, formalized at the end-pair
boundary level.} The geometric core of Step~4 is in
\path{TNLean/PEPS/TorusWindowExtraction.lean}: the reference edge $e$ is the
unique lattice edge joining the two end windows
(\path{isCrossingEdge_horizontalStaircaseEndWindows}, lifting the abstract
\path{isCrossingEdge_horizontalStaircase} to the actual
\path{horizontalStaircaseLeftWindow} and \path{horizontalStaircaseRightWindow}),
hence a boundary edge of each window
(\path{isRegionBoundaryEdge_horizontalStaircaseLeftWindow_referenceEdge} and
its right counterpart) but \emph{not} of the end pair
$S = W_0 \sqcup W_{L+K-1}$
(\path{not_isRegionBoundaryEdge_horizontalStaircaseEndPair_referenceEdge}):
both endpoints of $e$ lie in $S$, one in each window, so $e$ is the single
interior bond of $S$. Every other boundary edge of either window is a boundary
edge of $S$
(\path{isRegionBoundaryEdge_endPair_of_leftWindow} and its right counterpart),
so the boundary $\partial S$ is the disjoint union of the two windows' external
legs with $e$ the only shared interior bond. This is the geometric data the
bond-operator extraction consumes.

\medskip
\noindent\textbf{The matrix-intertwining route does not apply per window in
two dimensions.} The footnoted matrix-level form
\leanid{MPSChainTensor.exists_bondOperator_of_intertwine_span} extracts $X$
from spanning families $W_1, W_2$ of \emph{square} matrices
$\mathrm{Mat}_{D}$ with $F\,a \cdot W_2\,b = W_1\,a \cdot G\,b$. In the
one-dimensional chain it applies directly because a length-$L$ window has two
virtual legs --- a left bond and a right bond, both of dimension $D$ --- so the
window product $\text{\path{arcEval}}\,A\,r\,(\mathrm{ofFn}\,a)$ and the deformed
window $C\,0\,a$ are themselves $D \times D$ matrices. In two dimensions a
single end window touches the bond $e$ with \emph{one} endpoint, so the
$e$-leg is a single $\mathrm{Fin}(\mathrm{bondDim}\,e)$ index, not a pair: a
window block read at fixed external legs and physical configuration is a
\emph{vector} on $\mathrm{Fin}(\mathrm{bondDim}\,e)$, not a square matrix. The
two indices of the extracted $X(i,j)$ are the $e$-leg of $W_0$ (one) and the
$e$-leg of $W_{L+K-1}$ (the other); neither window alone supplies a square
matrix family, and folding a window's external perimeter into a second
$\mathrm{Fin}(\mathrm{bondDim}\,e)$ index would require
$\prod_{\text{external legs}} \mathrm{bondDim} = \mathrm{bondDim}\,e$, false in
general. So \path{exists_bondOperator_of_intertwine_span} cannot be fed from
per-window blocks.

The faithful two-dimensional extraction is the single-crossing
algebra-isomorphism route the torus Theorem~3 development already uses
(\path{exists_regionEdgeGauge_of_blockingData} of
\path{TNLean/PEPS/CoherentFrameInstance2.lean}): the region-inserted
coefficient \path{regionInsertedCoeff} inserts a matrix on the single boundary
edge $e$ and contracts the red block against the host, and Skolem--Noether
(\path{edgeGaugeFromInsertionAlgebraIsomorphism}) turns the resulting algebra
isomorphism into conjugation by an invertible edge matrix $Z$. The Theorem~3
engine inverts the \emph{host} $\mathrm{univ} \setminus \mathrm{red}$ (its
hypothesis \path{hCB}), which the corollary's minimal size does not provide; the
overlapping-window replacement plays the second end window $W_{L+K-1}$ in the
host's role, inverting only the two injective windows. Reaching $X$ from the
end-pair equality \path{staircasePair_insert_eq_open} therefore needs the
$\text{\path{extendInsert}}\,(W_0 \subseteq S)$-to-single-bond-$e$ peeling: reading the
corner extension on $S$ as a $\text{\path{regionInsertedCoeff}}$-style contraction of
$C_0$ against the genuine $W_{L+K-1}$ block across $e$. That peeling is the
host-boundary-edge embedding and the $Q$-weight span lemma recorded as the
unbuilt residual at the close of the previous section; the single-crossing
geometry above is the prerequisite it consumes, and is now in place.

\subsection{Step 5: the rest of the proof}

With the per-edge extraction in place, the remainder is the proof of
Theorem~3 specified by Ref.~\cite{Molnar2018NormalPEPS}, and the torus formalization
already realizes that part as a chain of layers: the insertion
correspondence turns the extraction into an algebra isomorphism of the
bond algebra at one reference edge per orientation, inner by
Skolem--Noether, giving a reference gauge with a coefficient identity;
translation transport spreads the reference witness to every edge as a
translation-covariant absorbed family; a comparison of a region $R$
with $R \cup \{v\}$ extracts the per-vertex scalar $\lambda$; and one
global contraction forces
$\lambda^{nm}=1$.\footnote{\leanid{TNLean.PEPS.exists_torusCovariantAbsorbedGaugeFamily},
\leanid{TNLean.PEPS.component_eq_gaugeVertex_of_cornerProportional}, and
\leanid{TNLean.PEPS.lambda_pow_card_torus_eq_one}, assembled in
\path{TNLean/PEPS/TorusFundamentalTheorem2.lean}.} Two of these layers
are anchored to the Theorem~3 shapes and must be re-anchored:

\begin{itemize}
  \item The comparison pair: the formalized pair is the $3\times3$
    square minus its corner against the full square. The
    one-orientation replacement is the $(L+1)\times(K+1)$ rectangle
    minus one corner vertex against the full $(L+1)\times(K+1)$
    rectangle; the smaller region is the union of the three $L\times K$,
    $L\times(K+1)$-, and $(L+1)\times K$-shaped unions of $L \times K$
    rectangles, the full rectangle is such a union as well, and both
    complements decompose into bands and one rectangle at
    $n \geq 2L+1$, $m \geq 2K+1$.
  \item The uniqueness clause of the corollary needs no system size and
    is already delivered by the formalized gauge-uniqueness statement
    for the torus family.
\end{itemize}

\section{What the sketch does not say, and what it does not need}

\subsection{No coherence question between rows}

A natural shortcut would treat each row of the torus as one site of a
coarse closed chain and apply the one-dimensional overlapping-window
theorem twice, row-wise and column-wise. That route produces one gauge
per column \emph{cut} --- an invertible matrix on the whole product of
the bond spaces crossing the cut --- and would then face a genuine
coherence problem: decomposing the cut gauges into one matrix per edge,
uniformly across rows, which the one-dimensional argument does not
provide. The sketch in Ref.~\cite{Molnar2018NormalPEPS} does not take this route and the question
never arises. The window chain is intrinsically two-dimensional and
local to one edge: the windows turn the corner around the edge (sliding
in both directions), and the end pair's single-crossing geometry forces
the extracted operator onto that edge alone. Coherence \emph{across}
edges is inherited from translation transport of one reference witness
per orientation, exactly as in Theorem~3.

The coherence problem is not merely unaddressed by the sketch; it is a
genuine obstruction, and the row-and-column shortcut cannot reach the
per-edge conclusion. The one-dimensional reduction reads a contracted row
only as an abstract family of matrices on the collected vertical-bond space
together with the closed-chain trace coefficients, and both inputs are
invariant under conjugating that family by an arbitrary invertible matrix
of the collected bond space. A non-product conjugation produces a contracted
row with the same coefficients and the same injectivity as a per-edge-gauged
one, so the reduction cannot tell them apart, while only the latter has a
per-edge factorization; and because an injective contracted row has scalar
self-conjugations, the reduction's gauge is that conjugator up to a scalar,
non-factoring whenever the conjugator is. The companion column-wise
application does not help: a row cut crosses the vertical edges of one row
boundary and a column cut the horizontal edges of one column boundary, two
disjoint edge sets sharing no bond, so the two cut gauges impose no mutual
consistency. This is recorded as a machine-checked refutation in
\path{TNLean/PEPS/TorusRowColumnReductionObstruction.lean}
(\leanid{TNLean.PEPS.RowColumnReductionObstruction.rowCutGaugeFactorizes_false}):
an injective contracted-row family and its non-product conjugate share their
state and injectivity, yet no per-edge product conjugator relates them.

\subsection{Why the formalized three-block route cannot reach the
corollary's sizes}

The formalized witness production for Theorem~3 blocks the torus into
three injective regions around the reference edge --- a red block, a
blue block meeting it only at the edge, and the complementary host ---
and the host must be injective. At the corollary's minimal size this
fails. Take $L = K = 2$, the $5\times5$ torus with columns and rows
$0,\dots,4$, the edge $e$ joining $(1,2)$ and $(2,2)$, and the natural
single-crossing pair: red $[0,1]\times[1,2]$, blue $[2,3]\times[2,3]$.
The host contains the vertex $(4,2)$, but every $2\times2$ square
containing $(4,2)$ --- including those wrapping the seam --- meets red
or blue. The host is therefore not a union of injective $2\times2$
regions, and its injectivity is not derivable from the hypotheses. The
overlapping-window route avoids the host entirely: the only inverted
complements are those of single windows and of consecutive-window
unions, all of which decompose at $(2L+1)\times(2K+1)$.

\subsection{Scope restrictions}

\begin{itemize}
  \item \textbf{One orientation suffices.} The corollary assumes
    injectivity of $L \times K$ regions only, with no transposed
    $K \times L$ hypothesis. The filled-in derivation confirms this:
    every region ever inverted is a union of $L \times K$ translates.
    The formalized Theorem~3 hypotheses
    (\leanid{TNLean.PEPS.NormalTorusRectangleInjectivityHypotheses})
    require both orientations, so the corollary needs a new
    one-orientation hypothesis structure.
  \item \textbf{The window must straddle the edge: $L, K \geq 2$.} The
    chain around a horizontal edge passes through windows containing
    both endpoints, which requires $L \geq 2$; vertical edges require
    $K \geq 2$. Ref.~\cite{Molnar2018NormalPEPS} states the corollary for unqualified $L, K$
    but proves the sketch at $L = K = 2$ and never addresses $L = 1$ or
    $K = 1$, where no window straddles the relevantly oriented edge and
    the chain cannot cross it. The formalization should carry
    $2 \leq L$ and $2 \leq K$ as explicit hypotheses; this is a
    locally-fixable clarification of the source statement, recorded
    here.
  \item \textbf{Region conventions.} As in the landed torus theorem,
    positive bond dimensions are assumed (the union lemma for
    overlapping regions needs them), and the formal conclusion is the
    torus surface of the source display: a translation-covariant
    per-edge family in place of the literal single pair $(X, Y)$, since
    the ordered-edge convention stores wraparound edges with swapped
    endpoints (see the section ``Closure on the torus'' of
    \path{docs/paper-gaps/peps_normal_ft_section3_route.tex}).
\end{itemize}

\section{Formalization plan}

The geometry, state-equality, and forward fixed-edge-operation layers of the
derivation above are complete. The remaining source step is the converse at
\path{Papers/1804.04964/paper_normal.tex:2368--2444} of
Ref.~\cite{Molnar2018NormalPEPS}: from physical operations on the overlapping
windows that
produce one state, recover the virtual operation $X$ on the highlighted edge
and prove the Lemma~5 maps from $O_1$ and $O_3^{\mathsf T}$ are algebra
homomorphisms.  The layer list below records the campaign around that step:

\begin{enumerate}
  \item \emph{Window geometry} (no tensors): the one-orientation
    rectangle-injectivity hypotheses; injectivity of all
    $a \times b$ rectangles with $a \geq L$, $b \geq K$ by sliding
    unions; the window family around an edge; disjointness of the end
    pair and the single-crossing characterization; the complement
    decompositions of single windows and consecutive-window unions at
    $n \geq 2L+1$, $m \geq 2K+1$, including the seam-wrapping bands.
  \item \emph{Deformed-window states}: the vocabulary for a torus state
    with one region's block replaced by an arbitrary tensor, and the
    consecutive-window comparison by complement inversion (the existing
    kernel-descent and two-block machinery applies once the complement
    injectivity is supplied by layer 1).
  \item \emph{Chaining and corner stripping}: extension of an
    open-boundary equality by contraction, and the completion-inversion
    step removing the corner block.
  \item \emph{The Lemma~5 converse on the staircase}: compare the first and
    last physical operations after adjoining the genuine corner tensors,
    invert the resulting regions, recover the unique virtual operation $X$,
    and prove the $O_1\mapsto X$ and $O_3^{\mathsf T}\mapsto X$
    homomorphisms by composition.
  \item \emph{The per-edge witness}: the insertion correspondence and
    Skolem--Noether at one reference edge per orientation, producing
    the same coefficient-identity witness shape the landed translation
    layers consume.
  \item \emph{Re-anchoring the downstream}: the covariant family, the
    $(L+1)\times(K+1)$ comparison pair, the scalar extraction, and the
    capstone with hypotheses: one-orientation $L \times K$ injectivity
    for both tensors, $2 \leq L$, $2 \leq K$, $n \geq 2L+1$,
    $m \geq 2K+1$, matched bond dimensions, positive bonds, the same
    state.
\end{enumerate}

Layer 1 is foundational and self-contained; layers 2--5 replace the
landed coarse three-site witness engine; layer 6 is a re-anchoring of
existing arguments. Until the capstone lands, the corollary of Ref.~\cite{Molnar2018NormalPEPS} at
\path{Papers/1804.04964/paper_normal.tex:2297--2318} remains
unformalized, and no blueprint entry should claim it.

Layer 1 is complete. The seam-avoiding window geometry --- the
one-orientation hypotheses, the general $a \times b$ injectivity by
sliding unions, the staircase end pair, and the single-crossing
characterization --- is in \path{TNLean/PEPS/TorusWindowRegion.lean}. The
seam-wrapping part --- the cyclic rectangle, the translation-invariant
window hypotheses, the consecutive-window unions, and the complement
decompositions of single windows and consecutive-window unions at
$n \geq 2L+1$, $m \geq 2K+1$ --- is in
\path{TNLean/PEPS/TorusWindowComplement.lean}. The seam-wrapping band is
modelled as a cyclic rectangle read through $\mathbb{Z}/n$ value
distance, so the full-height complement band of $\mathrm{width} - L$
columns is a single coordinate-translate of a rectangle and stays
injective by translation invariance of the tensor; the wraparound-free
restriction recovers the coordinate rectangle. The complement of each
single window and consecutive-window union is delivered both at the
abstract region-injectivity level and as
\path{RegionBlockedTensorInjective} of the tensor, the consumption shape
layer 2 needs.

The complement-inversion core of layer 2 is complete. The deformed-window
state --- the closed-torus contraction with one region's block replaced by
an arbitrary insert carrying that region's boundary structure --- is
\path{deformedRegionState}, defined through the block/complement pairing
of \path{sum_regionBlockedWeight_mul_complement}; taking the insert to be
the genuine block recovers the interior-bond multiple of the closed state
coefficient. The reusable engine
\path{deformedRegionState_insert_eq_of_complementInjective} strips a
closed-torus deformed-state equality to an open-boundary equality of the
two inserts on the region, once the set complement is blocked-tensor
injective: at a fixed region physical configuration the deformed state is
the complement tensor map applied to the insert, which is injective. Both
are in \path{TNLean/PEPS/TorusDeformedWindow.lean}. The engine is
specialized to the consecutive-window unions there as
\path{horizontalUnion_insert_eq_of_deformedState_eq} and its vertical
counterpart, the union complement injectivity supplied by layer~1. The
remaining content of
(\ref{eq:peps_normal_ft_2d_overlap_consecutive_window_open_equality}) ---
writing each single-window deformed
state as the union deformed state of the corner-extended insert
$C_j \sqcup T_{W_{j+1} \setminus W_j}$ --- is the contraction-extension of
the chaining layer (item~3 of the plan).

The geometry of the chaining and corner-stripping layer is in place. The
staircase patch $P$, its two end windows, the corner block, and the
completed corner rectangle are in
\path{TNLean/PEPS/TorusWindowChain.lean}, with: the disjointness of the two
end windows at $2L \le n$; the injectivity of each end window and of the
completed corner rectangle (each an $L \times K$ cyclic window); the patch
decomposition $P = S \sqcup \text{corner}$ at the no-wraparound bounds
$2L \le n$, $2K \le m$; and its corollary $P \setminus S = \text{corner}$,
the single block Step~3 completes and inverts. The decomposition reduces every
shifted cyclic distance to the base distance from the patch corner through the
shift arithmetic \path{zmod_val_sub_shift}.

The tensor content of the chaining and stripping is in
\path{TNLean/PEPS/TorusWindowChain2.lean}. The contraction-extension
\path{deformedRegionState_extend} writes a sub-region deformed state, read
as a function of the full physical configuration, as the deformed state of
a superset $S \supseteq R$ with the window's insert extended by the genuine
block on the added vertices $S \setminus R$. It is built from the disjoint
factorization of the complement weight $T_{\mathrm{univ} \setminus R}$ over
$\mathrm{univ} \setminus R = (S \setminus R) \sqcup
(\mathrm{univ} \setminus S)$: with the geometry $\mathrm{red} = R$,
$\mathrm{blue} = S \setminus R$,
$\mathrm{complement} = \mathrm{univ} \setminus S$, the landed
\path{ThreeBlockGeometry.regionInteriorBondProd_smul_regionBlockedWeight_threeBlockComplPhysical}
splits $T_{\mathrm{univ} \setminus R}$ into the blue-coupling combination of
the $T_{\mathrm{univ} \setminus S}$ weights, the sum over the
$\partial(\mathrm{univ} \setminus S)$ boundary configurations being exactly
the boundary-configuration split. The corner-extended insert
\path{extendInsert} contracts the window insert against the blue coupling
\path{ThreeBlockGeometry.threeBlockBlueCoeff}, scaled by the inverse of the
$\mathrm{univ} \setminus S$ interior-bond multiplicity the factorization
introduces (a nonzero scalar at positive bond dimensions, which cancels).
The deformed state is read region-independently through \path{assembleRegionσ},
so the extension equalities chain by transitivity. The corner extension thus
needs no standalone $\partial R \to \partial S$ boundary-configuration
equivalence; the three-block blue coupling supplies the split already
packaged.

With \path{deformedRegionState_extend}, the formal version of
(\ref{eq:peps_normal_ft_2d_overlap_consecutive_window_open_equality}),
\path{horizontalConsecutiveWindow_extend_eq}, writes each single-window state
as the union state of its corner-extended insert and passes the equality to
the consecutive-window engine, leaving an open-boundary equality of inserts
on each consecutive-window union. The patch chaining
\path{horizontalStaircase_patch_extend_eq} extends the two end windows'
states to the patch and chains them at the \emph{closed-torus} level, giving
an equality of the two patch assembled states. The faithful Step~3 stripping
\path{staircasePair_insert_eq_open} instead chains the patch at the
open-boundary level and cancels the shared injective completed corner $Q$,
leaving the open-boundary equality on the end pair $S$ with no
$\mathrm{univ} \setminus S$ hypothesis (the open-boundary route below).

The route to that stripping is the open-boundary chaining, not a complement
assembly. A closed-torus stripping would take the blocked-tensor injectivity
of $\mathrm{univ} \setminus S$ (the torus complement of the staircase end
pair) as an explicit hypothesis. As Step~3 records, that hypothesis is
\emph{not} derivable at the corollary's minimal
size: $\mathrm{univ} \setminus S$ is not a union of injective window
translates, because the end pair occupies all columns of one row and the
complement band there has width below $L$. Completing the corner does not
repair this: enlarging the compared region from $S$ to the patch (or to the
patch with the corner completed) leaves the closed-torus state unchanged, so
inverting the closed state on any superset $P' \supseteq S$ reduces to
inverting $\mathrm{univ} \setminus S$ rather than the added block $P'
\setminus S$. This collapse is recorded as
\path{deformedRegionStateAssembled_extendInsert_eq} and its inversion form
\path{deformedRegionStateAssembled_extendInsert_eq_of_subComplementInjective}
in \path{TNLean/PEPS/TorusWindowChain2.lean}: the superset closed-state route
consumes $\mathrm{univ} \setminus S$ injectivity and never the completed
corner's. The shared-corner cancellation of Step~3 is an open-boundary
operation, cancelling the injective completed rectangle
\path{horizontalStaircaseCompletedCorner} from an equality of inserts on the
completed region as tensors. The remaining work is therefore to re-derive the
patch chaining of Step~2 at the open-boundary level --- chaining the
consecutive-window open-boundary equalities
\path{horizontalConsecutiveWindow_extend_eq} across the patch by contraction
and transitivity, rather than collapsing to the patch closed state --- after
which the completed corner cancels locally by its injectivity alone.

The open-boundary route factors into two engines, neither of which inverts
$\mathrm{univ} \setminus S$. They are stated here with the machinery each needs.

\emph{The corner-extension composition.} For nested regions $R \subseteq S \subseteq P$ and an insert $C$ on
$R$, open-boundary chaining uses
\begin{align}
  \text{\path{extendInsert}} (S \subseteq P)\,
  \bigl(\text{\path{extendInsert}} (R \subseteq S)\, C\bigr)
    &= \text{\path{extendInsert}} (R \subseteq P)\, C
    &&\text{as inserts on } P.
  \label{eq:peps_normal_ft_2d_overlap_corner_extension_transitivity}
\end{align}
Each instance of
(\ref{eq:peps_normal_ft_2d_overlap_consecutive_window_open_equality}) lives
on the union $U_j$;
extending both sides through $\text{\path{extendInsert}}(U_j \subseteq P)$ and
collapsing the composition by
(\ref{eq:peps_normal_ft_2d_overlap_corner_extension_transitivity}) chains the equalities into one
patch-level insert equality. Equation
(\ref{eq:peps_normal_ft_2d_overlap_corner_extension_transitivity}) is
\emph{not} reducible to the
assembled-state collapse \path{deformedRegionState_extend}: both sides have the
same assembled state on $P$, but equal assembled states determine equal inserts
only when $\mathrm{univ} \setminus P$ is injective, which fails at the minimal
size. Stripping the inverse interior-bond multiplicities (the lemma
\path{extendInsert_const_smul} commutes the rescaling through), the identity
(\ref{eq:peps_normal_ft_2d_overlap_corner_extension_transitivity}) reduces to
the \emph{bare} composition law
\begin{align}
  \text{\path{bareExtendInsert}} (S \subseteq P)\,
  \bigl(\text{\path{bareExtendInsert}} (R \subseteq S)\, C\bigr)
    &= \mathrm{ribp}(\mathrm{univ} \setminus S)\cdot
       \text{\path{bareExtendInsert}} (R \subseteq P)\, C.
  \label{eq:peps_normal_ft_2d_overlap_bare_corner_extension_composition}
\end{align}
Equation~(\ref{eq:peps_normal_ft_2d_overlap_bare_corner_extension_composition})
is the blue-coupling composition of \path{ThreeBlockGeometry.threeBlockBlueCoeff}
across the two nested geometries, glued along the $\mathrm{univ} \setminus S$
cut. The product $\prod_{P \setminus R}$ splits as $\prod_{S \setminus R}
\cdot \prod_{P \setminus S}$, and the intermediate sum over
$\partial S$ glues the two global-configuration constrained sums along their
agreeing $\mathrm{univ} \setminus S$ boundary, with fiber multiplicity
$\mathrm{ribp}(\mathrm{univ} \setminus S)$. This is a four-block ($R$,
$S \setminus R$, $P \setminus S$, $\mathrm{univ} \setminus P$) merge-and-count,
beyond the existing three-block fiber tooling of
\path{TNLean/PEPS/RegionBlock/UnionInjectivityGeneral.lean}. It is landed in
\path{TNLean/PEPS/TorusWindowChain4.lean}: the merge-and-count is the two-sub-block
collapse \path{mergeCollapse2} over $\mathrm{univ} \setminus S$, with the two host
labels carried through by the boundary-edge embeddings
\path{regionBoundaryLabel_regionMerge_compl_of_subset} (the $\mathrm{univ}\setminus R$
host) and \path{regionBoundaryLabel_regionMerge_of_subset_left} (the
$\mathrm{univ}\setminus P$ host) and the product split \path{regionProd_split}; the
blue-coupling composition is \path{threeBlockBlueCoeff_comp}, the bare law
\path{bareExtendInsert_trans}, and the transitivity \path{extendInsert_trans}.

\emph{The shared-corner cancellation.} With the patch equality
$\text{\path{extendInsert}} (W_0 \subseteq P)\, C_0 = \text{\path{extendInsert}} (W_{L+K-1}
\subseteq P)\, C_{L+K-1}$ in hand, the corner cancellation is the injectivity of
$\text{\path{extendInsert}} (S \subseteq R)$ on inserts on $S$, where $R = S \sqcup Q$
and $Q = \text{\path{horizontalStaircaseCompletedCorner}}$ is the injective completed
rectangle: extending the patch equality to $R$ (one more transitivity step
through $P \subseteq R$, since $R = P \sqcup \text{added row}$) and cancelling
the shared $Q$-block gives the end-pair equality. The cancellation needs
\emph{only} $Q$ injective, never $\mathrm{univ} \setminus S$. For the nested
geometry $\mathrm{red} = S$, $\mathrm{blue} = R \setminus S = Q$,
$\mathrm{complement} = \mathrm{univ} \setminus R$, the corner-extended insert
contracts the $S$-insert against $\text{\path{ThreeBlockGeometry.threeBlockBlueCoeff}}$;
its $\sigma_{\mathrm{blue}}$-dependence is a combination of the $Q$ blocked
weights, so $Q$ injectivity (the chosen left inverse
\path{regionBlockedTensorMap_injective_of_injective} of the $Q$ block) reads off
the coupling rows, and the host-boundary decomposition $\partial(\mathrm{univ}
\setminus S) \hookrightarrow \partial Q \times \partial(\mathrm{univ}\setminus R)$
extracts the $S$-insert coefficients. This is the dual of the blue-strip
inversion \path{ThreeBlockGeometry.complCoeff_combination_eq_zero}: that lemma
pushes a host annihilation through the blue block to a complement-coupling
annihilation using $\mathrm{blue}$ injectivity; the cancellation here inverts in
the host index using $Q$ injectivity alone, the weaker hypothesis the union
lemma's two-step does not isolate.

The shared-corner cancellation is formalized in
\path{TNLean/PEPS/TorusWindowChain5.lean}. Since the corner extension is linear
in its insert, the injectivity of $\text{\path{extendInsert}}(S\subseteq R)$ reduces
to its kernel: the only insert on $S$ whose corner extension on $R$ vanishes is
the zero insert. The linearity lemmas \path{extendInsert_add},
\path{extendInsert_sub}, and \path{extendInsert_const_smul} supply this
reduction as \path{extendInsert_injective_of_kernel_trivial}. The physical
assembly of $R=S\sqcup Q$ is carried by \path{assembleSubRegionσ} and its two
restriction identities
\path{restrictSubRegionσ_assembleSubRegionσ} and
\path{restrictSubRegionσ_sdiff_assembleSubRegionσ}; these say that the $S$-
and $Q$-legs of a physical configuration on $R$ can be fixed independently.
The host-boundary extraction is
\path{ThreeBlockGeometry.regionBoundaryLabel_host_eq_hostLabelFrom}, which
identifies the host residual from the $Q$ and $\mathrm{univ}\setminus R$
residuals.

Together with the $Q$-weight span
\path{ThreeBlockGeometry.crossingBondProd_smul_threeBlockBlueCoeff_eq}, these
ingredients prove \path{extendInsert_kernel_trivial_of_addedInjective}: a
kernel vector is read through the $S\sqcup Q$ assembly, the
host-residual-weighted blue coupling vanishes for every $Q$-leg and complement
boundary, the blue block is stripped using $Q$ injectivity, and the
$S$-insert coefficients are recovered by host-label surjectivity at positive
bond dimensions. The cancellation theorem
\path{extendInsert_cancel_addedInjective} then cancels the injective added
block from an equality of corner extensions, using only that block's
injectivity and never the injectivity of $\mathrm{univ}\setminus S$. Its
staircase specialization \path{staircasePair_cancel} completes the end pair to
$S\sqcup Q$ with $Q$ the completed corner and cancels $Q$ from an
open-boundary equality of the two end-window inserts on the completed region.

The required open-boundary patch equality is formalized in
\path{TNLean/PEPS/TorusWindowChain6.lean}. The descending-arm comparison
\path{verticalConsecutiveWindow_extend_eq} is the transpose of the horizontal
comparison \path{horizontalConsecutiveWindow_extend_eq}: both give
open-boundary equalities of corner-extended inserts on a consecutive-window
union. The per-step theorem \path{staircaseStep_patch_extend_eq} lifts such a
consecutive-window equality from the union $U_j$ to the full patch $P$ by
extending both sides through $\text{\path{extendInsert}}(U_j\subseteq P)$ and using
\path{extendInsert_trans} along $W_j\subseteq U_j\subseteq P$. The chained
equality \path{staircasePatch_insert_eq} advances by induction from $W_0$ to
$W_{L+K-1}$; the intermediate windows disappear because their deformed states
are all equal to the common state.

Finally, \path{staircasePair_insert_eq_open} extends the patch equality through
the completed union $S\sqcup Q$, collapses the nested extensions by
\path{extendInsert_trans}, and applies \path{staircasePair_cancel}. The output
is the open-boundary equality on the end pair $S$. Its hypotheses are the
one-orientation window injectivity, the union closure, positive bonds,
$2\le L$, $2\le K$, $2L+1\le\mathrm{width}$, $2K+1\le\mathrm{height}$, and
the all-windows-agree input. No injectivity of $\mathrm{univ}\setminus S$ is
used. The obstruction that led to the Step~3 gap is therefore eliminated.

\section{The single-bond extraction and the host-injectivity scope of the
back half}

This section records two findings that follow the landing of the open-boundary
end-pair equality \path{staircasePair_insert_eq_open} and the single-crossing
geometry of Step~4: the geometric peeling of the end-pair equality onto the
single bond $e$, and a precise scope verdict on what host injectivity the
back-half of Theorem~3 (Step~5) actually needs.

\subsection{The boundary of the end pair is the windows' external legs}

Step~3 now delivers the open-boundary equality on the end pair $S = W_0 \sqcup
W_{L+K-1}$ without inverting $\mathrm{univ}\setminus S$: the corner-extension
composition
(\ref{eq:peps_normal_ft_2d_overlap_corner_extension_transitivity}) and the
shared-corner cancellation
$\text{\path{staircasePair\_cancel}}$ are landed, and the end-pair equality
\begin{align}
  \text{\path{extendInsert}}(W_0\subseteq S)\,C_0
    &= \text{\path{extendInsert}}(W_{L+K-1}\subseteq S)\,C_{L+K-1}
    &&\text{as inserts on } S.
  \label{eq:peps_normal_ft_2d_overlap_formal_end_pair_insert_equality}
\end{align}
is \path{staircasePair_insert_eq_open}. Equation
(\ref{eq:peps_normal_ft_2d_overlap_formal_end_pair_insert_equality}) has no
$\mathrm{univ}\setminus S$
hypothesis. The single-crossing geometry of Step~4 shows $e$ is the only lattice
edge joining the two end windows, hence the only interior bond of $S$.

The geometric peeling of
(\ref{eq:peps_normal_ft_2d_overlap_formal_end_pair_insert_equality}) onto the
bond $e$ rests on the complete
characterization of the boundary $\partial S$. The landed direction
($\text{\path{isRegionBoundaryEdge\_endPair\_of\_leftWindow}}$ and its right-window
counterpart) shows a boundary edge $f\neq e$ of one window is a boundary edge of
$S$. The converse, supplying the partition, is
$\text{\path{isRegionBoundaryEdge\_window\_of\_endPair}}$: a boundary edge of $S$ is a
boundary edge of one of the two windows, read off the in-region endpoint, whose
out-of-region endpoint lies outside that window since it lies outside all of $S$.
Combined with $e$'s interiority, this gives
$\text{\path{isRegionBoundaryEdge\_endPair\_iff}}$: $f$ is a boundary edge of $S$ iff it
is a boundary edge of one window and $f\neq e$; and
$\text{\path{isRegionBoundaryEdge\_endPair\_exactlyOne\_window}}$: each boundary edge of
$S$ lies on \emph{exactly} one window, since the only edge on both windows is $e$
($\text{\path{eq\_referenceEdge\_of\_isRegionBoundaryEdge\_both\_endWindows}}$), which is
not a boundary edge of $S$. So $\partial S$ is the disjoint union of the two
windows' external legs, with $e$ the single interior bond summed over in the
end-pair contraction. This is the geometric content the bond-$e$ coefficient
identity reads: the open legs of $S$ split into the left-window external legs and
the right-window external legs, and the matrix on $e$ is the only datum the two
windows share. These facts are in \path{TNLean/PEPS/TorusWindowPeeling/BoundaryGeometry.lean}.

\subsection{The back half transfers with a window as witness region}

The torus back half of Theorem~3 consumes host injectivity in exactly two
places, and at the corollary's minimal size $(2L+1)\times(2K+1)$ both are
supplied by one-orientation $L\times K$ window translates --- the obstruction of
Step~3 (that $\mathrm{univ}\setminus S$ is not injective at the minimal size) does
not propagate to either.

\emph{The witness region.} The per-edge interface
$\text{\path{EdgeCoeffIdentityWitness}}$ bundles a region whose boundary contains $e$,
the second tensor's region and host blocked-tensor injectivity, and the
conjugation coefficient identities on $e$. Its host field
$\text{\path{hCB}}:\text{\path{RegionBlockedTensorInjective}}\,B\,(\mathrm{univ}\setminus
\text{\path{region}})$ is host injectivity of the \emph{witness} region, not of $S$. The
rectangle route takes $\text{\path{region}}$ to be the reference rectangle's red block,
host injective only at the three-block sizes ($\geq 7$, or the seam-touching
$+5=\mathrm{width}$ special case). At the window route, $\text{\path{region}}$ may be a
single end window $W$: $e$ is a boundary edge of each end window
($\text{\path{isRegionBoundaryEdge\_horizontalStaircaseLeftWindow\_referenceEdge}}$ and
its right counterpart), $W$ is region injective
($\text{\path{horizontalStaircaseLeftWindow\_injective}}$), and crucially its host
$\mathrm{univ}\setminus W$ \emph{is} injective at $(2L+1)\times(2K+1)$ --- this is
the landed $\text{\path{regionBlockedTensorInjective\_windowComplement}}$, whose
complement decomposition (a full-height band of $\mathrm{width}-L\geq L+1$ columns
and an $L$-column band of $\mathrm{height}-K\geq K+1$ rows, each an $L\times K$
window union) needs exactly $2L+1\leq\mathrm{width}$, $2K+1\leq\mathrm{height}$.
So a single window is a valid $\text{\path{EdgeCoeffIdentityWitness}}$ region with host
injectivity available at the minimal size, once the peeling of \S5.1 produces the
bond-$e$ coefficient identity $\text{\path{regionInsertedCoeff}}\,A\,W\,\langle
e\rangle\,M = \text{\path{regionInsertedCoeff}}\,B\,W\,\langle e\rangle\,(Z\,M\,Z^{-1})$
with $W$ in place of the rectangle red block.

\emph{The comparison regions.} The scalar extraction
$\text{\path{lambda\_pow\_card\_torus\_eq\_one}}$ takes a single comparison region $R$
with $R$ and $\mathrm{univ}\setminus R$ injective; the per-vertex relation
$\text{\path{component\_eq\_gaugeVertex\_of\_cornerProportional}}$ compares $R$ against
$\text{\path{insert}}\,v\,R$, needing both regions' hosts injective. The formalized $R$
is $\text{\path{cornerRegion}}$ (the $3\times3$ square minus its corner) and
$\text{\path{insert}}\,v\,R = \text{\path{cornerSquare}}$ (the $3\times3$ square), at offset
$(2,2)$ with hardcoded $7\leq\mathrm{width}$, $7\leq\mathrm{height}$; their host
decompositions ($\text{\path{compl\_cornerRegion\_injective}}$,
$\text{\path{compl\_cornerSquare\_injective}}$) use both rectangle orientations
($\text{\path{rect23}}$, $\text{\path{rect32}}$, $\text{\path{wideRectangle}}$, $\text{\path{shortRectangle}}$)
and width-$2$/width-$(\mathrm{width}-5)$ bands tied to the $3\times3$ shape. These
are the $L=K=2$ instance of the source's $(L+1)\times(K+1)$ comparison pair. The
one-orientation replacement, as Step~5 of this note prescribes, is the
$(L+1)\times(K+1)$ cyclic rectangle minus one corner vertex against the full
$(L+1)\times(K+1)$ rectangle. Both are region injective with one orientation by
$\text{\path{arcRectangle\_injective}}$ (a cyclic rectangle of width $\geq L$, height
$\geq K$ is the union of the $L\times K$ windows sliding over it). Both hosts
decompose at the minimal size: by $\text{\path{compl\_torusArcRectangle}}$ the complement
of a $(L+1)\times(K+1)$ cyclic rectangle is a full-height band of
$\mathrm{width}-(L+1)\geq L$ columns and an $(L+1)$-column band of
$\mathrm{height}-(K+1)\geq K$ rows, each injective by
$\text{\path{arcRectangle\_injective}}$; the band widths reach exactly $L$ and the
heights exactly $K$ at $2L+1$, $2K+1$, which is why the corollary needs no larger
size. The minus-corner region differs by one vertex and decomposes the same way
with one extra corner rectangle, the transpose of the formalized
$\text{\path{cornerRegion}}$ band decomposition.

\emph{Verdict.} The back half transfers to the corollary's minimal size with the
window as witness region and the $(L+1)\times(K+1)$-shaped comparison pair, both
host-decomposable into one-orientation $L\times K$ window translates at $2L+1$,
$2K+1$. No host injectivity genuinely fails at the minimal size for any region the
back half inverts: the sole failure, $\mathrm{univ}\setminus S$, is confined to
the Step~3 stripping, which the open-boundary route already circumvents. The
capstone is therefore an \emph{assembly} of landed machinery --- the peeling of
\S5.1 to the bond-$e$ coefficient identity (the one genuinely-new mechanical
piece, whose geometric foundation is now landed), the window-region witness with
its already-available host injectivity, and the re-anchored comparison pair built
from $\text{\path{arcRectangle\_injective}}$ and $\text{\path{compl\_torusArcRectangle}}$ ---
not further research. A host-free variant of the back-half engines is \emph{not}
needed.

\subsection{The bond-insertion bridge and the single-bond peeling are landed}

The peeling of \S5.1 onto the single bond $e$ rests on a bridge between the two
vocabularies the back half mixes: the \emph{deformed-window} inserts the
overlapping-window staircase produces, and the \emph{region-inserted coefficient}
$\text{\path{regionInsertedCoeff}}$ the per-edge gauge of Step~5 reads.  These were two
parallel developments with no connecting lemma.  The bridge is now in
\path{TNLean/PEPS/TorusWindowPeeling/BondInsertedRegion.lean}.

The bond-inserted insert $\text{\path{bondInsertedRegionInsert}}\,A\,R\,f\,M$ on a region
$R$ is the insert whose boundary configuration $\nu$ (read by the complement
contraction) carries the $M$-coupled combination of the genuine $R$ block over the
boundary configurations agreeing with $\nu$ away from $f$.  Its deformed state is
exactly the region-inserted coefficient:
$\text{\path{deformedRegionState\_bondInsertedRegionInsert}}$ proves
$\text{\path{deformedRegionState}}\,A\,R\,(\text{\path{bondInsertedRegionInsert}}\,A\,R\,f\,M) =
\text{\path{regionInsertedCoeff}}\,A\,R\,f\,M$, the inner boundary sum of the insert
reproducing the double boundary sum of the coefficient.  Composing with the
corner-extension identity $\text{\path{deformedRegionState\_extend}}$ gives
$\text{\path{regionInsertedCoeff\_eq\_extendInsert\_bondInserted}}$: the coefficient on $R$,
read off a global configuration, equals the assembled deformed state on a superset
$S$ of the corner-extended bond-inserted insert.  With $R$ an end window and $S$ the
end pair, the genuine block of the \emph{opposite} window the bond $e$ joins to is
carried by the corner extension --- the geometric content the
$\mathrm{univ}\setminus W = W' \sqcup (\mathrm{univ}\setminus S)$ split of \S5.1
describes, now realized through the landed corner extension rather than a bespoke
contraction split.

The single-bond peeling itself is
$\text{\path{regionInsertedCoeff\_endWindows\_eq\_of\_staircase}}$ in
\path{TNLean/PEPS/TorusWindowPeeling/SingleBond.lean}.  Feeding the end-pair equality
$\text{\path{staircasePair\_insert\_eq\_open}}$ the bond-inserted inserts on the two end
windows (and arbitrary inserts on the intermediate windows), all sharing one deformed
state, and bridging both sides through
$\text{\path{regionInsertedCoeff\_eq\_extendInsert\_bondInserted}}$, peels the equality onto
$e$: the two end windows' region-inserted coefficients with $M$ and $M'$ on $e$ agree,
read off any global configuration.  The reference edge is a boundary edge of each end
window by the landed single-crossing geometry
($\text{\path{isRegionBoundaryEdge\_staircaseWindow\_zero\_referenceEdge}}$ and its last-window
counterpart), so $M$ and $M'$ live on the same bond.

The window-region witness packaging of \S5.2 is
$\text{\path{windowEdgeCoeffIdentityWitness}}$ (and its hypotheses form
$\text{\path{windowEdgeCoeffIdentityWitness\_of\_hypotheses}}$) in
\path{TNLean/PEPS/TorusWindowWitness.lean}: given the bond-$e$ conjugation coefficient
identity, every other $\text{\path{EdgeCoeffIdentityWitness}}$ field with the left end window
as region is discharged from the landed window geometry, the host injectivity supplied
at the minimal size by $\text{\path{regionBlockedTensorInjective\_windowComplement}}$.  The
witness is the reference-edge interface the torus covariant-family machinery consumes.

\medskip
\noindent\textbf{The residual.}  What remains of \S5.1 is the production of the
bond-$e$ \emph{conjugation} coefficient identity the witness packaging consumes ---
$\text{\path{regionInsertedCoeff}}\,A\,W\,\langle e\rangle\,M =
\text{\path{regionInsertedCoeff}}\,B\,W\,\langle e\rangle\,(Z\,M\,Z^{-1})$.  The Skolem--Noether
back end of this production is now landed and reduces it to a single cross-tensor input.
The reusable region-level algebra of $\text{\path{RegionInsertionTransfer}}$ derives the gauge
$Z$ and the conjugation form of the coefficient identity from a \emph{coefficient
transfer} on the window: the read-off $\text{\path{exists\_regionEdgeGauge\_of\_coeffTransfer}}$
turns the two cross-tensor coefficient transfers and the forward-transfer multiplicativity
into $Z$ and the conjugation identity (the lemma
$\text{\path{exists\_regionConjCoeffIdentity\_of\_coeffTransfer}}$), and feeding that to the
window-region witness packaging discharges the witness's conjugation field (the lemma
$\text{\path{exists\_windowEdgeCoeffIdentityWitness\_of\_coeffTransfer}}$).  These are in
\path{TNLean/PEPS/TorusWindowMult.lean}.  All region-level algebra --- additivity,
homogeneity, identity preservation, the two-sided inverse, the Skolem--Noether read-off
--- is unconditional on the window and host injectivity, both available at the minimal
size, so the window witness assembles from the coefficient transfer alone, with no
conjugation field assumed.

What remains is therefore exactly the staircase \emph{coefficient transfer} on the
window: for each inserted matrix $M$ on $A$'s edge $e$ a matching matrix $N$ on $B$ with
equal window coefficient
($\text{\path{htransferAB}}$), its reverse ($\text{\path{htransferBA}}$), and the multiplicativity of
the chosen forward transfer ($\text{\path{hmul}}$).  The same-tensor family that feeds the
single-bond peeling is now complete.  The first window carries the boundary insert of
$M^{\mathsf T}$, the windows $W_1,\ldots,W_{L-1}$ carry the interior-bond insert of $M$,
and the remaining windows carry the boundary insert of $M$.  Their common state is the
closed torus state with $M$ inserted on the same edge $e$, with the common
non-boundary-bond factor supplied by translation invariance.  The construction and its
end-window identity are
$\text{\path{staircaseVirtualOperationInsert}}$ and
$\text{\path{regionInsertedCoeff\_endWindows\_eq\_staircaseVirtualOperation}}$.

The residual is the cross-tensor multiplicative transfer: the homomorphism the rectangle route
reads from the physical realization of the inserted operator.  The region-level recovery
$\text{\path{regionInsertionTransfer\_of\_realizes}}$ derives all three transfer fields ---
$\text{\path{htransferAB}}$, $\text{\path{htransferBA}}$, and $\text{\path{hmul}}$ --- from one region
physical-to-virtual realization $\text{\path{RegionTransferRealizes}}$ in each direction, with
no further hypothesis; so the window route's residual collapses to producing that
realization on the window from the staircase, in place of the edge-middle injectivity
($\mathrm{univ}\setminus\{u,v\}$ injective) the rectangle route uses, which fails at the
minimal size --- the same obstruction, at the edge, that $\mathrm{univ}\setminus S$
exhibits at the end pair.  A coefficient transfer alone does not yield $\text{\path{hmul}}$: the
multiplicativity is the homomorphism of the physical insertion operator, which the
single-bond geometry localizes to one edge but does not on its own make multiplicative.

\section{The multiplicativity verdict: it is the coarse three-site
super-bond, not a single-window span}

This section records the definitive resolution of the one residual question the
overlapping-window corollary reduces to: whether the forward-transfer
multiplicativity $\text{\path{hmul}}$ at the left end window is obtainable from window and
host injectivity at the corollary's minimal size, or genuinely needs a block-level
construction the landed single-window machinery does not provide.

\subsection{The single-window span does not exist, and is not the mechanism}

The natural reading of $\text{\path{hmul}}$ as a single-window fact is the question whether
the reference-edge-restricted matrix family of the left end window $W$ --- the bond
matrix on $e$ read off the window block at fixed external legs and physical
configuration --- spans the full $\mathrm{Fin}(\mathrm{bondDim}\,e)$ algebra, the way
a one-dimensional length-$L$ window product spans the full $D\times D$ algebra and
feeds the homomorphism extraction $\text{\path{exists\_insertionHom}}$ through
$\text{\path{LinearMap.ext\_on\_range}}$. It does not. A one-dimensional window has exactly
two virtual legs, both of dimension $D$, so its product \emph{is} a square matrix on
the bond and window injectivity is bond-algebra spanning directly. A two-dimensional
$L\times K$ end window touches $e$ with \emph{one} endpoint, so its $e$-restricted
family is a family of \emph{vectors} on $\mathrm{Fin}(\mathrm{bondDim}\,e)$, and the
window's whole perimeter is open. Folding the perimeter into a second
$\mathrm{Fin}(\mathrm{bondDim}\,e)$ index would require
$\prod_{\text{external legs}}\mathrm{bondDim} = \mathrm{bondDim}\,e$, false in
general; the full boundary-configuration family is linearly independent (window block
injectivity), but its projection onto the $e$-leg at fixed external legs is not a
square-matrix span. So the one-dimensional template port has no two-dimensional
single-window analogue, and the homomorphism cannot be read from one window's block.

\subsection{The faithful mechanism is the coarse super-bond, available without
single-vertex injectivity, but obstructed for the single-window red block at the
minimal size}

Ref.~\cite{Molnar2018NormalPEPS} supplies $\text{\path{hmul}}$ not from a single window but from the homomorphism
of the physical insertion operator, which the formalization realizes through the
\emph{coarse three-site} route already landed for the rectangle blocking: three
blocked regions (a red block, a blue block meeting it only at $e$, and the
complement) are assembled into a coarse tensor on a three-vertex graph, whose three
super-sites are injective directly from the three blocked-region injectivities, with
no single-vertex injectivity of the original tensor. The coarse tensor is
edge-blocked three-site injective at the distinguished super-bond, so the
single-vertex transfer-matrix multiplicativity $\text{\path{edgeTransferMatrix\_mul}}$
applies \emph{at the super-bond}; conjugated by the two bond models this is
$\text{\path{regionEdgeTransfer\_mul}}$, and through the identification of the chosen forward
transfer with the concrete one ($\text{\path{coeffTransferMap\_eq\_regionEdgeTransfer}}$, by
$\text{\path{regionInsertedCoeff\_injective}}$) it is exactly $\text{\path{hmul}}$ for the chosen
transfer over the red block. The super-vertex injectivity stands in for the
single-vertex spanning the vertex-injective route would use; this is why the route
needs only blocked-region injectivity of the three regions.

This mechanism applies to the window route iff the left end window $W$ can play the
coarse red block, which requires a blue block making $e$ the single red-to-blue
crossing and a complement $\mathrm{univ}\setminus(W\cup\mathrm{blue})$ injective. The
single-crossing requirement is met canonically: the right end window $W'$ is the only
block joined to $W$ by exactly the reference edge $e$
($\text{\path{isCrossingEdge\_horizontalStaircaseEndWindows}}$). But that forces
$W\cup\mathrm{blue} = W\cup W' = S$, the staircase end pair, so the complement is
$\mathrm{univ}\setminus S$ --- the very region Step~3 proves is \emph{not} injective
at the corollary's minimal size, because $S$ occupies all columns of one row and the
complement band there has width below $L$. So the coarse three-site frame with
$\mathrm{red}=W$ and the natural single-crossing partner is unbuildable at the minimal
size: its $\text{\path{complement\_injective}}$ field is exactly the obstructed
$\mathrm{univ}\setminus S$ injectivity, the same obstruction the open-boundary
staircase route was built to circumvent at the \emph{state-equality} level, now
reappearing at the \emph{multiplicativity} level. Enlarging the blue block beyond
$W'$ cannot repair this within the single-window red frame: a blue block making $e$
the sole red-to-blue crossing must contain $e$'s out-of-$W$ endpoint and avoid
creating any second crossing of $\partial W$, and any such enlargement only shrinks
$\mathrm{univ}\setminus(W\cup\mathrm{blue})$, making its injectivity no easier than for
$\mathrm{univ}\setminus S$. The faithful repair is not a larger partner but a
different mechanism, below.

\subsection{Verdict}

The verdict is \textbf{(b)}: the $e$-bond multiplicativity at the minimal size is
\emph{not} derivable from window and host injectivity of the single left end window.
It is the coarse three-site super-bond multiplicativity, which needs a third injective
block --- the complement of the red-blue pair --- and at the minimal size, with the
red block a single window and the single-crossing partner forced to be the opposite
window, that complement is $\mathrm{univ}\setminus S$, non-injective. The read-off
side ($\text{\path{htransferAB}}$, $\text{\path{htransferBA}}$, and the whole region-level gauge
algebra) \emph{is} powered by window and single-window-complement injectivity, both
available at the minimal size; only the homomorphism is not.

The overlapping-window realization sketched in Ref.~\cite{Molnar2018NormalPEPS} keeps the same virtual operation
$M$ on $e$ while representing it on every staircase window.  That common-state family
is now built below.  It closes the single-operation part of this discussion but not the
cross-tensor homomorphism: the remaining task is to compare the two tensor families and
show that the resulting coefficient transfer respects products, following Lemma~5's
composition of the maps from $O_1$ and $O_3^{\mathsf T}$ rather than a single-window
span or an assumed coarse frame.

\subsection{The foundation that lands}

The reusable multiplicativity is exposed as
\leanid{TNLean.PEPS.coeffTransferMap\_mul\_of\_coherentFrames}
(\path{TNLean/PEPS/RegionBlock/CoarseThreeSiteMul.lean}): two coherent coarse blocking
frames over the two tensors sharing the three regions, the bond dimensions, and the
single-crossing edge make the chosen forward transfer over the red block
multiplicative, with no single-vertex injectivity. This is the $\text{\path{hmul}}$ field
extracted from the bundled gauge assembly
\leanid{TNLean.PEPS.exists\_regionEdgeGauge\_of\_coherentFrames}, so a consumer holding
only the coefficient transfer can obtain the multiplicativity from a frame. It makes
the verdict auditable in Lean: the window route's $\text{\path{hmul}}$ is one application of
this theorem with $\mathrm{red}=W$, and the unmet hypothesis is exactly the frame's
$\text{\path{complement\_injective}}$ at $\mathrm{univ}\setminus S$, the documented
obstruction.  This theorem remains available for any future coherent frame, but the
fixed-edge common-state family below does not require such a frame.  The remaining work
is the cross-tensor multiplicativity argument.

\subsection{The common-state family is complete}

The earlier analysis incorrectly claimed that the intermediate staircase windows contain
neither endpoint of the reference edge and therefore require transport across different
bonds.  The machine-checked endpoint classification gives the opposite conclusion:
$W_1,\ldots,W_{L-1}$ contain both endpoints of $e$, while
$W_L,\ldots,W_{L+K-1}$ contain its ordered left endpoint. Every window in
the family of Ref.~\cite{Molnar2018NormalPEPS} therefore contains at least
one endpoint of the same distinguished edge $e$.

The missing family is now
\leanid{TNLean.PEPS.staircaseVirtualOperationInsert}
(\path{TNLean/PEPS/TorusWindowVirtualOperationFamily.lean}).  It uses the boundary insert
of $M^{\mathsf T}$ on $W_0$, the interior-bond insert of $M$ on
$W_1,\ldots,W_{L-1}$, and the boundary insert of $M$ on the remaining windows.  The
interior construction
\leanid{TNLean.PEPS.interiorBondInsertedRegionInsert}
(\path{TNLean/PEPS/RegionBlock/InteriorBondInsertion.lean}) starts from the ordered left
endpoint and adjoins the genuine tensors of the window by corner extension.  This formalizes the virtual insert asserted at
\path{Papers/1804.04964/paper_normal.tex:2320--2321} of
Ref.~\cite{Molnar2018NormalPEPS} and uses the same adjoining contraction as
the end-region comparison at
\path{Papers/1804.04964/paper_normal.tex:2415--2444} of that reference.  The physical operator itself is
\leanid{TNLean.PEPS.staircaseVirtualOperationPhysicalOp}: it is the prescribed insert
linear combination composed with a left inverse of the injective window block.  The theorem
\leanid{TNLean.PEPS.staircaseVirtualOperationPhysicalOp\_realizes} proves the open-boundary statement of
Ref.~\cite{Molnar2018NormalPEPS}
$O_j(M)T_{B,W_j}(\mu)=C_j(M)(\mu)$ for every~$\mu$, before complementary contraction.
The insert's assembled state is the window's non-boundary bond product times
$\mathrm{edgeInsertedCoeff}(e,M)$.  Translation
invariance makes that product uniform over the staircase, proving one common deformed
state for all $L+K$ inserts.

Consequently
\leanid{TNLean.PEPS.regionInsertedCoeff\_endWindows\_eq\_staircaseVirtualOperation}
supplies the source-faithful single-tensor identity between the first-window insertion of
$M^{\mathsf T}$ and the last-window insertion of $M$.  This closes the common-state-family
residual.  This end-window coefficient identity follows directly from the two region-to-edge
identities and translation invariance; it needs neither complement inversion nor a peeling
hypothesis.  It does not by itself construct the cross-tensor transfer from $A$ to $B$ or
prove that transfer multiplicative; those are the remaining Lemma~5-style comparison
steps.

\section{Normality and discarded stronger routes}

Ref.~\cite{Molnar2018NormalPEPS} calls a PEPS normal when blocking
suitable regions produces an injective tensor.\footnote{\path{Papers/1804.04964/paper_normal.tex:1318}.}
In the corollary under consideration, every $L\times K$ rectangle is assumed
injective, so those rectangles themselves witness normality.  No
single-vertex injectivity is asserted.  This distinction is essential:
single-vertex injectivity is strictly stronger than injectivity after
blocking, and therefore cannot be inserted into a formal statement of the
corollary.

The formal window hypothesis has exactly this content: every cyclic
$L\times K$ window is injective as a map from its virtual boundary space to
its physical space.\footnote{Formal declarations:
\leanid{TNLean.PEPS.NormalTorusArcWindowInjectivityHypotheses} and
\leanid{TNLean.PEPS.RegionBlockedTensorInjective}.}  It neither assumes
single-vertex injectivity nor omits a further normality condition from the
source.

Two stronger mechanisms available elsewhere in the development are therefore
only conditional alternatives.  The first reads a bond operation from a
single vertex and requires single-vertex injectivity.  The second proves
multiplicativity from a coherent three-region blocking frame.\footnote{Formal
declaration:
\leanid{TNLean.PEPS.coeffTransferMap\_mul\_of\_coherentFrames}.}  Neither
hypothesis occurs in the corollary, and neither is needed for the forward
fixed-edge family proved above.

\subsection{Why the different-bond auxiliary route is irrelevant}

Ref.~\cite{Molnar2018NormalPEPS} keeps one highlighted edge $e$ fixed.  The first staircase window
contains its ordered right endpoint, the next $L-1$ windows contain both
endpoints, and the remaining windows contain its ordered left endpoint.
Consequently, the same matrix $X$ can be represented on every window; there
is no step in which an operation must be transported from one boundary bond
to a different boundary bond of the same window.

Such a transport would in any case require information not contained in
window injectivity.  A two-dimensional window is injective with respect to
its whole multi-bond perimeter.  After all other boundary legs are fixed, its
restriction to two selected bonds need not span a square matrix algebra.  The
one-dimensional overlapping-window argument has no such difficulty because
a chain window has exactly two virtual boundary legs, so its blocked tensors
are already square matrices on the bond.

The natural three-region substitute also lies outside the source route at the
minimal size.  Two $L\times K$ windows with a single crossing edge form the
diagonally offset staircase pair, whose union is not a cyclic rectangle.  On
the shared row the complement has fewer than $L$ consecutive columns, so the
window-union argument used elsewhere does not supply injectivity of the third
region of a coherent frame.\footnote{Formal geometry:
\leanid{TNLean.PEPS.horizontalStaircasePair\_ne\_arcRectangle} and
\leanid{TNLean.PEPS.horizontalStaircasePair\_complement\_band\_lt\_window}
in \path{TNLean/PEPS/TorusWindowSingleCrossingObstruction.lean}.}

These observations rule out two auxiliary shortcuts; they do not challenge
the paper's fixed-edge comparison.  The remaining source step is the converse at
\path{Papers/1804.04964/paper_normal.tex:2368--2444} of
Ref.~\cite{Molnar2018NormalPEPS}: compare the physical operations on the overlapping
windows, adjoin and invert the two corner regions, recover the unique virtual
operation $X$, and prove the maps $O_1\mapsto X$ and
$O_3^{\mathsf T}\mapsto X$ multiplicative by the composition argument of
Lemma~5.

\bibliographystyle{alpha}
\bibliography{references}

\end{document}
